% =====================================================================
% DROP-IN SECTION for futon5 holes/tech-notes/paper/main.tex
% Written by lon-claude-6, 2026-07-17, at Joe's request: "try to write up how
% they act on policies, because that was our big idea about what they might be
% useful for, and it isn't in the paper yet."
%
% NOT committed to main.tex directly: oxf-claude-2 is actively editing that file
% (343940a, 2022ebf), and my futon5 worktree is detached at an ancestor with a
% live job running from it. Handing over rather than conflicting.
%
% Uses the paper's existing macros \sig and \bit{}. Suggested placement: a new
% \section after \section{The Propagator} and before \section{The Fail Bank},
% since it depends on \ref{sec:fixedpoints} and \ref{sec:conjugacy} and nothing
% later.
%
% CONTAINS A CORRECTION TO \subsection{Exact fixed points} — see the flagged
% paragraph. That subsection contradicts the paper's own new theorem.
% =====================================================================

\section{Propagators Acting on Policies}
\label{sec:policies}

Everything above treats the propagator as an operator on a byte. This section
makes a claim we have not previously stated: \textbf{the propagator's carrier is
a policy, and it has been one all along}. The consequence is a second level at
which the whole apparatus applies --- and a negative result when we run it there.

\subsection{A policy is a colouring; so is the carrier}

The propagator is $g[\sig(k)] := \nu(g[k])$ over a colouring $g : N \to P$ with a
shape map $\sig : N \to N$ and a value map $\nu : P \to P$. A \emph{policy} is a
function from conditions to actions. These are the same type. Not by analogy:
$N$ is the condition set, $P$ is the action set, and $g$ is the prescription.

This is not a new construction. \textbf{The MetaCA rule is already a policy.} A
rule is a map $\mathrm{NBHD} \to \{0,1\}$ --- ``when you see neighbourhood $k$,
respond $\bit{k}$'' --- which is a prescription indexed by a condition. And
\texttt{rule-permute} (generator.clj:732) reads, literally,
\[
  \bit{\sig(k)} \;:=\; \lnot\,\bit{k}, \qquad k \sim \mathrm{Uniform}\{0..7\}.
\]
That is $g[\sig(k)] := \nu(g[k])$ with $g$ the rule's truth table, $\sig$ the
operator, and $\nu = \mathrm{Bool.not}$. \textbf{The genotype's truth table is an
eight-cell system whose topology is $\sig$: a CA inside the CA.} This is why the
Lean development is about \texttt{Bool.not} specifically --- it is not a
convenient toy, it is the operator this paper runs.

Three objects in this stack are policies, and we had not noticed they were the
same object:
\begin{enumerate}
  \item \textbf{The rule}, $\mathrm{NBHD} \to \{0,1\}$. Level 0.
  \item \textbf{The switch}, $\{\textsf{bored}, \textsf{active}\} \to \{\sig,
        \textsc{no-op}\}$ (Section~\ref{sec:baldwin}). A two-condition policy over
        propagators.
  \item \textbf{A pattern cascade}, $\mathrm{BINS} \to \mathrm{ACTS}$, where
        $\mathrm{ACTS}$ are themselves propagators. Level 1.
\end{enumerate}
The third is the interesting one, because there $P$ is the very thing this paper
studies: a propagator rewriting a policy whose values are propagators is the
theory applied to its own controller.

\subsection{The multiplication rule}

Given a cascade $g : \mathrm{BINS} \to \mathrm{ACTS}$ and an arm $a$, what is
$g \cdot a$? Not $\nu = a$: an arm is an \emph{element} of $\mathrm{ACTS}$, not a
map on it. But arms are permutations of the eight neighbourhoods, so they compose,
and each arm induces $\nu_a(\sig) = a \circ \sig$ --- left multiplication, the
regular representation of $S_8$ on itself. Cayley then gives, before any run:
\[
  \nu_a \text{ has a fixed point} \iff a = \mathrm{id}.
\]
So every $g \cdot a$ is fixed-point-free except $g \cdot \mathrm{id}$. Two readings
follow, and they fail in opposite ways:
\begin{center}
\begin{tabular}{llll}
\toprule
reading & rule & $\mathrm{FREE}$ & outcome \\
\midrule
pointwise  & $g' = a \circ g$              & $\emptyset$ & always exists; \emph{no scaffold} \\
propagator & $g'[\sig(k)] = a \circ g(k)$  & $\{\textsf{rich}\}$ & \emph{no solution} \\
\bottomrule
\end{tabular}
\end{center}

\paragraph{The propagator reading is unsatisfiable.}
With the clamped shift $\sig(k) = \max(k-1,0)$ transported to the regime scale,
$\sig(\textsf{dead}) = \textsf{dead}$ is a self-loop, whose constraint reads
$g[\textsf{dead}] = \nu(g[\textsf{dead}])$ --- a fixed point of $\nu$, which Cayley
has just forbidden. Confirmed by exhaustive search over all 40{,}320 candidate
roots: \textsc{unsat} for every $a \neq \mathrm{id}$. The unique satisfiable case
is $a = \mathrm{id}$, whose witness is the \emph{constant colouring} --- this
paper's own death configuration, arriving as the only solution.

\paragraph{The scaffold is what kills it.}
Run with proper parallel-update semantics (simultaneous update, merged writes,
contradiction reported rather than resolved), the rewrite \textbf{settles to
period 1 in three steps for every $a$, including every fixed-point-free one}.
$\textsf{rich}$ is \textsc{free} --- no writer --- hence a \emph{constant source},
and $\sig$ is a shift register fed by it, so the settled policy is
$g_\infty[k] = a^{\mathrm{depth}(k)} \circ \sig_{\textsf{rich}}$.

\textbf{Settling here is caused by the shape, not by $\nu$.} This is the converse
side of the fixed-point theorem: a fixed point of $\nu$ is \emph{sufficient} for
settling, but fixed-point-\emph{freeness} is \emph{not sufficient to prevent} it.
And it inverts the intended reading of $\mathrm{FREE}$: the surjective shape
($\mathrm{FREE} = \emptyset$, no scaffold) cycles forever with period
$\mathrm{ord}(a)$, while the shape \emph{with} a scaffold settles every time. A
free cell is a cell with no input; a cell with no input is a constant; a constant
floods the chain. Life needs feedback, and this $\sig$'s functional graph has
exactly one cycle: the self-loop, of length one, which is the one demanding a
$\nu$ fixed point.

\subsection{We ran it. It does not work, and the reason is instructive}

We ran the pointwise products as controllers on the tokamak (12 seeds, identical
across arms). Nothing died. That is not a result: an identity controller also
dies 0/12 while doing nothing. The informative row is worse than a null.

\textbf{$g \cdot \sig_{5127}$ is vacuous.} $\sig_{5127}$ has order 2, so at depth 0
the product is $\sig \circ \sig = \mathrm{id}$: the policy plays the identity in
$\textsf{rich}$, never leaves $\textsf{rich}$, and therefore never consults its
other four boxes. Its spacetime diagram is \emph{pixel-identical} to the
do-nothing control (\texttt{compare -metric AE} $= 0$ over 243{,}936 pixels). A
five-condition policy in which one condition fires and that condition does nothing.

Every test we had passed it. The $\nu$ fixed-point test passed; the satisfiability
check passed; it died 0/12. \textbf{Nothing in the apparatus detects vacuity},
because vacuity is not death --- it is a controller with nothing to control. This
is the mirror of Section~\ref{sec:census}'s class-4 finding, where diversity 1
reads as maximally dead and is maximally selected. There the metric says dead and
the system computes; here the metric says alive and the system idles. Neither
diversity nor death rate measures what is happening.

\subsection{Why level 0 works and level 1 does not}

The two levels run the same operator and behave oppositely. The difference is not
in the propagator. It is that \textbf{at level 0 the propagator is interrupted,
and at level 1 it is not}.

Two independent results in this paper already say so, and we had not connected
them:
\begin{itemize}
  \item \textbf{Blending is not a detail} (Section~\ref{sec:fixedpoints}): the
        isolated attractors are low-weight bytes; the census's collapse targets are
        majority and identity. Blending overwrites the propagator's input before it
        can converge.
  \item \textbf{The no-op is the mechanism} (Section~\ref{sec:compositions}):
        either propagator alone collapses; inside
        $\mathrm{switch}(\textsf{bored}, \cdot, \textsc{no-op})$ it sustains
        diversity 43--63; \emph{two live branches give noise}. Holding is what lets
        structure persist.
\end{itemize}
Both are the same statement. \textbf{A propagator alone is not a generative
mechanism; it is half of one.} The other half is an interrupter --- something that
prevents it reaching the fixed point it is otherwise guaranteed to reach. Whether
the interrupter overwrites the input (blending) or skips the step (\textsc{no-op})
is secondary.

The level-1 cascade has \emph{no} interrupter. It settles in three steps. That is
the whole negative result, and it was predictable from the level-0 findings the
paper already contains.

\paragraph{A design flaw this exposes.}
The cascade has no hold branch: every one of its five conditions fires a live
propagator. Worse, the arm \emph{named} $\mathrm{identity}$ is not a no-op ---
$\sig = \mathrm{id}$ means $\bit{k} := \lnot\,\bit{k}$, a flip at every draw, which
is \emph{maximal} churn. \textbf{The quietest-sounding arm in the set is the
noisiest.} By Section~\ref{sec:compositions}'s finding, a policy all of whose
branches fire should give noise --- which is exactly what
$g \cdot \sig_{5127}$ became. The testable repair is to give the policy a genuine
\textsc{no-op} branch, which no cascade we have written has.

\paragraph{What blending would be at the policy level.}
Level 0's interrupter mixes a cell's rule with its neighbours'. The analogue at
level 1 is to mix a policy with its neighbours' policies --- that is, a
\emph{population} of cascades that blend, rather than one cascade rewritten in
isolation. We have not run this. It yields a sharp prediction: a single cascade
under a propagator settles to period 1 in three steps (measured, every arm),
whereas a population of cascades under a propagator \emph{and} a blend should not.
That is the experiment this section owes.

\subsection{A correction to \texorpdfstring{Section~\ref{sec:fixedpoints}}{Exact fixed points}}

\textbf{Flagged for the authors; not yet reconciled.} Section~\ref{sec:fixedpoints}
states that no operator has a single-byte attractor, that support $\le 1$ occurs
for 0.0\% of $\sig$ with an observed minimum support of 4, and gives as the reason
that a fixed byte ``would require $\bit{\sig(k)} = \lnot\,\bit{k}$ for all $k$
simultaneously, which an always-invert operator forbids.''

That reason is not right, and the paper's own newer theorem says so: a Boolean
fixed point exists \emph{iff every cycle of $\sig$ is even}. An always-invert
operator forbids a fixed byte only on \emph{odd} cycles. Since 27.34\% of $S_8$ is
all-even, 27.34\% of operators \emph{do} have fixed bytes.

Enumerating all 256 bytes against \texttt{gen/rule-permute} itself:
\begin{center}
\begin{tabular}{lll}
\toprule
$\sig$ & cycle type & fixed bytes \\
\midrule
$\mathrm{rotate}{+}1$ & $[8]$          & \textbf{2} \quad \texttt{01001101}, \texttt{10110010} \\
$\mathrm{rotate}{+}2$ & $[4,4]$        & \textbf{4} \quad \texttt{00101011}, \texttt{01100110}, \texttt{10011001}, \texttt{11010100} \\
$\mathrm{three{+}five}$ & $[3,5]$      & 0 \\
$\sig_{5127}$ & $[1,1,2,2,2]$          & 0 \\
$\mathrm{id}$ & $[1^8]$                & 0 \\
\bottomrule
\end{tabular}
\end{center}
Exactly $2^{k}$ for all-even $\sig$ with $k$ cycles, and $0$ for any-odd. The new
theorem predicts these counts precisely; Section~\ref{sec:fixedpoints} contradicts
them. We believe the subsection predates the theorem and is stale rather than
wrong-in-substance --- but as written the two cannot both stand, and the
``minimum support is 4'' figure needs re-deriving.

This also sharpens the ``blending is not a detail'' paragraph in the authors'
favour. \textbf{Isolated attainment is certain, not merely possible}: solving the
propagator exactly as a $256 \times 256$ chain over uniform draws gives
$P(\text{absorbed}) = 1.000$ for both $\mathrm{rotate}{+}1$ and
$\mathrm{rotate}{+}2$. So the gap between the theorem and the census is not
``exists but is not reached''. It is that \emph{the propagator never runs to
absorption}, because the wiring is
$\mathrm{rule} := \texttt{rule-permute}(\mathrm{blend}(\cdot))$ --- one step per
generation on a freshly blended input. \textbf{The propagator is a bias, not an
absorber.} The theorem describes a system MetaCA never runs in isolation, which is
precisely what Section~\ref{sec:fixedpoints} was reaching for.

\subsection{Disposition}

The claim is a type identity: policies are the propagator's carrier, and MetaCA's
rule was always one. The claim is cheap and we believe it. The \emph{application}
is not banked: run at level 1, the propagator settles, and the one product that
survived is vacuous. What the negative result buys is a mechanism --- a propagator
needs an interrupter, and the level-1 object has none --- and a named next
experiment: a population of policies that blend.
